%% Andrea Piccolo - CV tailored for Halliburton Landmark (Data Integration Geoscientist)
%%
%% Compile with: lualatex main_halliburton.tex
%% (pdflatex often fails on modern MiKTeX with fontawesome5 font-expansion errors.)

\documentclass[11pt,a4paper,sans]{moderncv}
\moderncvstyle{banking}
\moderncvcolor{blue}

% Force both first and last name AND section headings to render in moderncv
% blue (color1). Default banking on lualatex+MiKTeX leaves these black, which
% looks inconsistent with the rest of the blue accent scheme.
\renewcommand*{\firstnamestyle}[1]{{\fontsize{34}{36}\bfseries\upshape\color{color1}#1}}
\renewcommand*{\lastnamestyle}[1]{{\fontsize{34}{36}\bfseries\upshape\color{color1}#1}}
\renewcommand*{\sectionstyle}[1]{{\sectionfont\color{color1}#1}}

\usepackage[utf8]{inputenc}
\usepackage{hyperref}
\hypersetup{
    colorlinks=true,
    linkcolor=blue,
    filecolor=magenta,
    urlcolor=blue,
    pdftitle={Andrea Piccolo - CV},
    pdfpagemode=FullScreen,
}
\usepackage[scale=0.77]{geometry}
\usepackage{needspace}
\usepackage{import}

% personal data
\name{Andrea}{Piccolo}
\address{Leeds, UK}{}{}
\phone[mobile]{+44 7831 007579}
\email{andreapiccolope@gmail.com}
\extrainfo{\href{https://www.linkedin.com/in/ap0889/}{LinkedIn}, \href{https://github.com/APiccolo89}{GitHub} \textbar{} UK Global Talent Visa -- no sponsorship required}

\begin{document}

\makecvtitle

% ============================================================
%     PROFILE STATEMENT
% ============================================================

\vspace{6pt}
\small{Computational geoscientist with a PhD in Computational Geodynamics and a decade of experience combining geoscience domain knowledge with Python-based data engineering and numerical modelling. Condensed 30 TB of simulated geoscientific data into a curated, publicly accessible 300 GB database and optimised numerical software to a 10\texttimes{} speedup, working across research institutions in Germany and the UK.}

% ============================================================
%     SKILLS
% ============================================================

\section{Skills}
\vspace{1pt}
\begin{itemize}

\item \textbf{Geoscience Domain Expertise}: Stratigraphy, petrology, structural geology and geochemistry (MSc/BSc training); crustal mechanics and geomechanical modelling from a decade of computational geodynamics research

\item \textbf{Python \& ETL Pipelines}: NumPy, SciPy, pandas, h5py, Numba, pytest — building automated ETL workflows that clean, structure, and curate large geoscientific datasets for reliable downstream use

\item \textbf{Geomechanical Modelling}: Forward mechanical modelling contributed to an open-source basin modelling package (Julia), thermo-mechanical and finite element modelling (FEniCSx), physical process coupling within large-scale 3D numerical codes

\item \textbf{Data Engineering at Scale}: TB-scale dataset condensation and curation (30 TB to 300 GB public databases), 100+ TB simulation data processing, developing relational database and SQL skills

\item \textbf{Software Engineering \& Collaboration}: Git, GitHub Actions (CI/CD), regression testing, cross-institutional collaboration, communicating technical results to varied technical audiences

\end{itemize}

% ============================================================
%     PROFESSIONAL EXPERIENCE
% ============================================================

\section{Professional Experience}
\vspace{3pt}
\begin{itemize}

\needspace{5\baselineskip}
\item{\cventry{02/2025--Present}{Research Fellow}{University of Leeds}{Leeds, UK}{}{\vspace{1pt}
\begin{itemize}
    \item Developed a numerical software package using FEniCSx and Python; optimised a pre-existing package achieving a 10\texttimes{} speedup
    \item Built automated ETL workflows to run simulations, process outputs, and analyse large geoscientific datasets
    \item Integrated regression tests to validate numerical simulation results against benchmarks
    \item Managed the full analytical workflow independently, from model design to documentation and reporting
\end{itemize}}}

\vspace{3pt}

\needspace{5\baselineskip}
\item{\cventry{11/2021--05/2024}{Research Fellow}{University of Bayreuth}{Bayreuth, Germany}{}{\vspace{1pt}
\begin{itemize}
    \item Extended an open-source modelling package in Julia; merged forward mechanical modelling into basin modelling
    \item Delivered publicly available databases: 30 TB of simulated data condensed and curated into 300 GB of accessible outputs
    \item Developed Python ETL scripts for data cleaning, structuring, and visualisation to support project reporting
    \item Cross-institutional collaboration with the KIT Karlsruhe basin modelling group, translating partner data requirements into shared workflows
\end{itemize}}}

\vspace{3pt}

\needspace{5\baselineskip}
\item{\cventry{11/2015--07/2021}{Research Fellow}{Johannes Gutenberg Universität Mainz}{Mainz, Germany}{}{\vspace{1pt}
\begin{itemize}
    \item Incorporated chemical modelling into a large-scale 3D numerical code
    \item Generated and processed 100+ TB of simulation data; developed 2D/3D visualisation tools
    \item Presented results at international conferences; supported supervision and teaching activities
\end{itemize}}}

\end{itemize}

% ============================================================
%     EDUCATION
% ============================================================

\section{Education}
\vspace{1pt}
\begin{itemize}

\needspace{5\baselineskip}
\item{\cventry{11/2015--05/2019}{PhD in Computational Geodynamics}{Johannes Gutenberg Universität Mainz}{Mainz, Germany}{}{\vspace{1pt}
Thesis: ``Continental crust production during the early Earth: Constraints from geodynamic models.''
}}

\vspace{3pt}

\item{\cventry{10/2012--03/2015}{MSc in Geology and Technical Geology}{University of Padua}{Padua, Italy}{}{\vspace{1pt}
Data analysis (chemical, petrological, stratigraphical), numerical modelling, scientific communication, field data acquisition.
}}

\vspace{3pt}

\item{\cventry{10/2008--03/2012}{BSc in Earth Sciences}{University of Padua}{Padua, Italy}{}{\vspace{1pt}
Data analysis, MATLAB, field data acquisition.
}}

\end{itemize}

% ============================================================
%     LANGUAGES
% ============================================================

\section{Languages}
\vspace{1pt}
\begin{itemize}
\item Italian (native), English (C1), German (B1).
\end{itemize}

% ============================================================
%     REFERENCES
% ============================================================

\section{References}
\vspace{1pt}
\begin{itemize}
\item Prof.\ Dr.\ Timothy J.\ Craig --- University of Leeds (\href{mailto:T.J.Craig@leeds.ac.uk}{T.J.Craig@leeds.ac.uk})
\item Prof.\ Dr.\ Marcell Thielmann --- University of Bonn (\href{mailto:thielmann@geo.uni-bonn.de}{thielmann@geo.uni-bonn.de})
\item Prof.\ Dr.\ Nevena Andric Tomasevic --- KIT Karlsruhe (\href{mailto:nevena.tomasevic@kit.edu}{nevena.tomasevic@kit.edu})
\item Dr.\ Arne Spang --- University of Bayreuth (\href{mailto:arne.spang@uni-bayreuth.de}{arne.spang@uni-bayreuth.de})
\end{itemize}

\end{document}
